% Figure: which face state each dual face uses (Chapter 6). Cells, liquid,
% unknowns and every dual face of component 1 with its class, read from the
% code's schedules (single-phase schedule with its boundary dual faces, phase
% schedule with boundary unknowns, regularity flags, CTU switch):
% generated/momentum_rules.tex (momentum.jl, momentum_rules_data).
\begin{figure}
  \centering
  \def\fxrnone{%
    \def\fxrcell##1##2##3##4##5{}\def\fxrliquid##1{}\def\fxrplic##1##2##3##4{}%
    \def\fxrdof##1##2##3##4{}\def\fxrlink##1##2##3##4##5##6##7{}}%
  % Classes of \fxrlink: 0 regular (scheme), 1 non-regular and 2 incident to a
  % boundary unknown (amr_interface_scheme), 3 CTU, 4--6 boundary dual faces
  % (B1)--(B3). Never colour alone: the interface scheme is dashed, CTU thick,
  % the boundary dual faces end on the wall.
  \tikzset{
    fxr 0/.style={draw=black!65, line width=0.45pt},
    fxr 1/.style={draw=trred!85!black, line width=0.75pt, dash pattern=on 2.2pt off 1.2pt},
    fxr 2/.style={fxr 1},
    fxr 3/.style={draw=trgreen!65!black, line width=1.5pt},
    fxr 4/.style={draw=trblue, line width=1.0pt},
    fxr 5/.style={draw=trblue!70, line width=0.5pt, dash pattern=on 0.7pt off 0.9pt},
    fxr 6/.style={draw=trpurple, line width=1.2pt},
    fxr bend 0/.style={}, fxr bend 1/.style={bend left=40},
    fxr bend -1/.style={bend right=40},
    fxr dof/.style={circle, fill=black, inner sep=0pt, minimum size=2.6pt},
    fxr bdof/.style={rectangle, draw=black, fill=white, line width=0.5pt,
      inner sep=0pt, minimum size=3.4pt},
  }%
  % Panel #1: 1 single-phase convection, 2 two-phase transport.
  \def\fxrpanel#1{%
    \begin{tikzpicture}[x=1.45cm,y=1.45cm]
      \fxrnone
      \ifnum#1=1
        \def\fxrcell##1##2##3##4##5{\fill[tr level ##5] (##1,##2) rectangle (##3,##4);}%
      \else
        \def\fxrliquid##1{\fill[tr liquid] ##1;}%
      \fi
      % Generated by docs/techref/figures/generate.jl; do not edit.
\fxrcell{2}{0}{3}{1}{0}
\fxrliquid{(2,0) -- (3,0) -- (3,1) -- (2,1) -- cycle}
\fxrcell{3}{0}{4}{1}{0}
\fxrliquid{(3,0) -- (4,0) -- (4,1) -- (3,1) -- cycle}
\fxrcell{2}{1}{3}{2}{0}
\fxrliquid{(2,1) -- (3,1) -- (3,2) -- (2.571429,2) -- (2,1.8) -- cycle}
\fxrplic{2.571429}{2}{2}{1.8}
\fxrcell{3}{1}{4}{2}{0}
\fxrliquid{(3,1) -- (4,1) -- (4,2) -- (3,2) -- cycle}
\fxrcell{0}{2}{1}{3}{0}
\fxrcell{1}{2}{2}{3}{0}
\fxrcell{0}{3}{1}{4}{0}
\fxrcell{1}{3}{2}{4}{0}
\fxrcell{2}{2}{3}{3}{0}
\fxrliquid{(2.571429,2) -- (3,2) -- (3,2.15) -- cycle}
\fxrplic{2.571429}{2}{3}{2.15}
\fxrcell{3}{2}{4}{3}{0}
\fxrliquid{(3,2) -- (4,2) -- (4,2.5) -- (3,2.15) -- cycle}
\fxrplic{4}{2.5}{3}{2.15}
\fxrcell{2}{3}{3}{4}{0}
\fxrcell{3}{3}{4}{4}{0}
\fxrcell{0}{0}{0.5}{0.5}{1}
\fxrliquid{(0,0) -- (0.5,0) -- (0.5,0.5) -- (0,0.5) -- cycle}
\fxrcell{0.5}{0}{1}{0.5}{1}
\fxrliquid{(0.5,0) -- (1,0) -- (1,0.5) -- (0.5,0.5) -- cycle}
\fxrcell{0}{0.5}{0.5}{1}{1}
\fxrliquid{(0,0.5) -- (0.5,0.5) -- (0.5,1) -- (0,1) -- cycle}
\fxrcell{0.5}{0.5}{1}{1}{1}
\fxrliquid{(0.5,0.5) -- (1,0.5) -- (1,1) -- (0.5,1) -- cycle}
\fxrcell{1}{0}{1.5}{0.5}{1}
\fxrliquid{(1,0) -- (1.5,0) -- (1.5,0.5) -- (1,0.5) -- cycle}
\fxrcell{1.5}{0}{2}{0.5}{1}
\fxrliquid{(1.5,0) -- (2,0) -- (2,0.5) -- (1.5,0.5) -- cycle}
\fxrcell{1}{0.5}{1.5}{1}{1}
\fxrliquid{(1,0.5) -- (1.5,0.5) -- (1.5,1) -- (1,1) -- cycle}
\fxrcell{1.5}{0.5}{2}{1}{1}
\fxrliquid{(1.5,0.5) -- (2,0.5) -- (2,1) -- (1.5,1) -- cycle}
\fxrcell{0}{1}{0.5}{1.5}{1}
\fxrliquid{(0,1) -- (0.5,1) -- (0.5,1.275) -- (0,1.1) -- cycle}
\fxrplic{0.5}{1.275}{0}{1.1}
\fxrcell{0.5}{1}{1}{1.5}{1}
\fxrliquid{(0.5,1) -- (1,1) -- (1,1.45) -- (0.5,1.275) -- cycle}
\fxrplic{1}{1.45}{0.5}{1.275}
\fxrcell{0}{1.5}{0.5}{2}{1}
\fxrcell{0.5}{1.5}{1}{2}{1}
\fxrcell{1}{1}{1.5}{1.5}{1}
\fxrliquid{(1,1) -- (1.5,1) -- (1.5,1.5) -- (1.142857,1.5) -- (1,1.45) -- cycle}
\fxrplic{1.142857}{1.5}{1}{1.45}
\fxrcell{1.5}{1}{2}{1.5}{1}
\fxrliquid{(1.5,1) -- (2,1) -- (2,1.5) -- (1.5,1.5) -- cycle}
\fxrcell{1}{1.5}{1.5}{2}{1}
\fxrliquid{(1.142857,1.5) -- (1.5,1.5) -- (1.5,1.625) -- cycle}
\fxrplic{1.142857}{1.5}{1.5}{1.625}
\fxrcell{1.5}{1.5}{2}{2}{1}
\fxrliquid{(1.5,1.5) -- (2,1.5) -- (2,1.8) -- (1.5,1.625) -- cycle}
\fxrplic{2}{1.8}{1.5}{1.625}
\fxrlink{1}{0.5}{0.25}{0.5}{0.75}{0}{0}
\fxrlink{1}{0.5}{0.25}{1}{0.25}{0}{0}
\fxrlink{1}{0.5}{0.75}{0.5}{1.25}{0}{0}
\fxrlink{1}{0.5}{0.75}{1}{0.75}{0}{0}
\fxrlink{1}{0.5}{1.25}{0.5}{1.75}{0}{0}
\fxrlink{1}{0.5}{1.25}{1}{1.25}{0}{0}
\fxrlink{1}{0.5}{1.75}{1}{1.75}{0}{0}
\fxrlink{1}{0.5}{1.75}{1}{2.5}{1}{0}
\fxrlink{1}{1}{0.25}{1}{0.75}{0}{0}
\fxrlink{1}{1}{0.25}{1.5}{0.25}{0}{0}
\fxrlink{1}{1}{0.75}{1}{1.25}{0}{0}
\fxrlink{1}{1}{0.75}{1.5}{0.75}{0}{0}
\fxrlink{1}{1}{1.25}{1}{1.75}{0}{0}
\fxrlink{1}{1}{1.25}{1.5}{1.25}{0}{0}
\fxrlink{1}{1}{1.75}{1}{2.5}{1}{0}
\fxrlink{1}{1}{1.75}{1.5}{1.75}{0}{0}
\fxrlink{1}{1}{2.5}{1}{3.5}{0}{0}
\fxrlink{1}{1}{2.5}{1.5}{1.75}{1}{0}
\fxrlink{1}{1}{2.5}{2}{2.5}{0}{0}
\fxrlink{1}{1}{3.5}{2}{3.5}{0}{0}
\fxrlink{1}{1.5}{0.25}{1.5}{0.75}{0}{0}
\fxrlink{1}{1.5}{0.25}{2}{0.25}{1}{0}
\fxrlink{1}{1.5}{0.75}{1.5}{1.25}{0}{0}
\fxrlink{1}{1.5}{0.75}{2}{0.75}{1}{0}
\fxrlink{1}{1.5}{1.25}{1.5}{1.75}{0}{0}
\fxrlink{1}{1.5}{1.25}{2}{1.25}{1}{0}
\fxrlink{1}{1.5}{1.75}{2}{1.75}{1}{0}
\fxrlink{1}{1.5}{1.75}{2}{2.5}{1}{0}
\fxrlink{1}{2}{0.25}{2}{0.75}{1}{1}
\fxrlink{1}{2}{0.25}{2}{0.75}{1}{-1}
\fxrlink{1}{2}{0.25}{2}{1.25}{1}{0}
\fxrlink{1}{2}{0.25}{2}{1.75}{1}{0}
\fxrlink{1}{2}{0.25}{3}{0.5}{1}{0}
\fxrlink{1}{2}{0.75}{2}{1.25}{1}{0}
\fxrlink{1}{2}{0.75}{2}{1.75}{1}{0}
\fxrlink{1}{2}{0.75}{3}{0.5}{1}{0}
\fxrlink{1}{2}{1.25}{2}{1.75}{1}{1}
\fxrlink{1}{2}{1.25}{2}{1.75}{1}{-1}
\fxrlink{1}{2}{1.25}{2}{2.5}{1}{0}
\fxrlink{1}{2}{1.25}{3}{1.5}{1}{0}
\fxrlink{1}{2}{1.75}{2}{2.5}{1}{0}
\fxrlink{1}{2}{1.75}{3}{1.5}{1}{0}
\fxrlink{1}{2}{2.5}{2}{3.5}{0}{0}
\fxrlink{1}{2}{2.5}{3}{2.5}{0}{0}
\fxrlink{1}{2}{3.5}{3}{3.5}{0}{0}
\fxrlink{1}{3}{0.5}{3}{1.5}{0}{0}
\fxrlink{1}{3}{1.5}{3}{2.5}{0}{0}
\fxrlink{1}{3}{2.5}{3}{3.5}{0}{0}
\fxrlink{1}{0.5}{0.25}{0}{0.25}{4}{0}
\fxrlink{1}{0.5}{0.75}{0}{0.75}{4}{0}
\fxrlink{1}{0.5}{1.25}{0}{1.25}{4}{0}
\fxrlink{1}{0.5}{1.75}{0}{1.75}{4}{0}
\fxrlink{1}{1}{2.5}{0}{2.5}{4}{0}
\fxrlink{1}{1}{3.5}{0}{3.5}{4}{0}
\fxrlink{1}{3}{0.5}{4}{0.5}{4}{0}
\fxrlink{1}{3}{1.5}{4}{1.5}{4}{0}
\fxrlink{1}{3}{2.5}{4}{2.5}{4}{0}
\fxrlink{1}{3}{3.5}{4}{3.5}{4}{0}
\fxrlink{1}{0.5}{0.25}{0.25}{0}{5}{0}
\fxrlink{1}{0.5}{0.25}{0.75}{0}{5}{0}
\fxrlink{1}{1}{0.25}{0.75}{0}{5}{0}
\fxrlink{1}{1}{0.25}{1.25}{0}{5}{0}
\fxrlink{1}{1.5}{0.25}{1.25}{0}{5}{0}
\fxrlink{1}{1.5}{0.25}{1.75}{0}{5}{0}
\fxrlink{1}{2}{0.25}{1.75}{0}{5}{0}
\fxrlink{1}{2}{0.25}{2.5}{0}{5}{0}
\fxrlink{1}{2}{0.75}{2.5}{0}{5}{0}
\fxrlink{1}{3}{0.5}{2.5}{0}{5}{0}
\fxrlink{1}{3}{0.5}{3.5}{0}{5}{0}
\fxrlink{1}{1}{3.5}{0.5}{4}{5}{0}
\fxrlink{1}{1}{3.5}{1.5}{4}{5}{0}
\fxrlink{1}{2}{3.5}{1.5}{4}{5}{0}
\fxrlink{1}{2}{3.5}{2.5}{4}{5}{0}
\fxrlink{1}{3}{3.5}{2.5}{4}{5}{0}
\fxrlink{1}{3}{3.5}{3.5}{4}{5}{0}
\fxrlink{1}{0.5}{1.75}{0}{2.5}{6}{0}
\fxrdof{1}{0.5}{0.25}{0}
\fxrdof{1}{0.5}{0.75}{0}
\fxrdof{1}{0.5}{1.25}{0}
\fxrdof{1}{0.5}{1.75}{0}
\fxrdof{1}{1}{0.25}{0}
\fxrdof{1}{1}{0.75}{0}
\fxrdof{1}{1}{1.25}{0}
\fxrdof{1}{1}{1.75}{0}
\fxrdof{1}{1}{2.5}{0}
\fxrdof{1}{1}{3.5}{0}
\fxrdof{1}{1.5}{0.25}{0}
\fxrdof{1}{1.5}{0.75}{0}
\fxrdof{1}{1.5}{1.25}{0}
\fxrdof{1}{1.5}{1.75}{0}
\fxrdof{1}{2}{0.25}{0}
\fxrdof{1}{2}{0.75}{0}
\fxrdof{1}{2}{1.25}{0}
\fxrdof{1}{2}{1.75}{0}
\fxrdof{1}{2}{2.5}{0}
\fxrdof{1}{2}{3.5}{0}
\fxrdof{1}{3}{0.5}{0}
\fxrdof{1}{3}{1.5}{0}
\fxrdof{1}{3}{2.5}{0}
\fxrdof{1}{3}{3.5}{0}
\fxrlink{2}{0.5}{0.25}{0.5}{0.75}{0}{0}
\fxrlink{2}{0.5}{0.25}{1}{0.25}{0}{0}
\fxrlink{2}{0.5}{0.25}{0}{0.25}{2}{0}
\fxrlink{2}{0.5}{0.75}{0.5}{1.25}{3}{0}
\fxrlink{2}{0.5}{0.75}{1}{0.75}{0}{0}
\fxrlink{2}{0.5}{0.75}{0}{0.75}{2}{0}
\fxrlink{2}{0.5}{1.25}{0.5}{1.75}{3}{0}
\fxrlink{2}{0.5}{1.25}{1}{1.25}{3}{0}
\fxrlink{2}{0.5}{1.25}{0}{1.25}{3}{0}
\fxrlink{2}{0.5}{1.75}{1}{1.75}{3}{0}
\fxrlink{2}{0.5}{1.75}{1}{2.5}{1}{0}
\fxrlink{2}{0.5}{1.75}{0}{1.75}{2}{0}
\fxrlink{2}{0.5}{1.75}{0}{2.5}{2}{0}
\fxrlink{2}{1}{0.25}{1}{0.75}{0}{0}
\fxrlink{2}{1}{0.25}{1.5}{0.25}{0}{0}
\fxrlink{2}{1}{0.75}{1}{1.25}{3}{0}
\fxrlink{2}{1}{0.75}{1.5}{0.75}{0}{0}
\fxrlink{2}{1}{1.25}{1}{1.75}{3}{0}
\fxrlink{2}{1}{1.25}{1.5}{1.25}{3}{0}
\fxrlink{2}{1}{1.75}{1}{2.5}{3}{0}
\fxrlink{2}{1}{1.75}{1.5}{1.75}{3}{0}
\fxrlink{2}{1}{2.5}{1}{3.5}{0}{0}
\fxrlink{2}{1}{2.5}{1.5}{1.75}{3}{0}
\fxrlink{2}{1}{2.5}{2}{2.5}{3}{0}
\fxrlink{2}{1}{2.5}{0}{2.5}{2}{0}
\fxrlink{2}{1}{3.5}{2}{3.5}{0}{0}
\fxrlink{2}{1}{3.5}{0}{3.5}{2}{0}
\fxrlink{2}{1.5}{0.25}{1.5}{0.75}{0}{0}
\fxrlink{2}{1.5}{0.25}{2}{0.25}{1}{0}
\fxrlink{2}{1.5}{0.75}{1.5}{1.25}{3}{0}
\fxrlink{2}{1.5}{0.75}{2}{0.75}{1}{0}
\fxrlink{2}{1.5}{1.25}{1.5}{1.75}{3}{0}
\fxrlink{2}{1.5}{1.25}{2}{1.25}{3}{0}
\fxrlink{2}{1.5}{1.75}{2}{1.75}{3}{0}
\fxrlink{2}{1.5}{1.75}{2}{2.5}{3}{0}
\fxrlink{2}{2}{0.25}{2}{0.75}{1}{1}
\fxrlink{2}{2}{0.25}{2}{0.75}{1}{-1}
\fxrlink{2}{2}{0.25}{2}{1.25}{3}{0}
\fxrlink{2}{2}{0.25}{2}{1.75}{3}{0}
\fxrlink{2}{2}{0.25}{3}{0.5}{1}{0}
\fxrlink{2}{2}{0.75}{2}{1.25}{3}{0}
\fxrlink{2}{2}{0.75}{2}{1.75}{3}{0}
\fxrlink{2}{2}{0.75}{3}{0.5}{1}{0}
\fxrlink{2}{2}{1.25}{2}{1.75}{3}{1}
\fxrlink{2}{2}{1.25}{2}{1.75}{3}{-1}
\fxrlink{2}{2}{1.25}{2}{2.5}{3}{0}
\fxrlink{2}{2}{1.25}{3}{1.5}{3}{0}
\fxrlink{2}{2}{1.75}{2}{2.5}{3}{0}
\fxrlink{2}{2}{1.75}{3}{1.5}{3}{0}
\fxrlink{2}{2}{2.5}{2}{3.5}{3}{0}
\fxrlink{2}{2}{2.5}{3}{2.5}{3}{0}
\fxrlink{2}{2}{3.5}{3}{3.5}{0}{0}
\fxrlink{2}{3}{0.5}{3}{1.5}{3}{0}
\fxrlink{2}{3}{0.5}{4}{0.5}{2}{0}
\fxrlink{2}{3}{1.5}{3}{2.5}{3}{0}
\fxrlink{2}{3}{1.5}{4}{1.5}{3}{0}
\fxrlink{2}{3}{2.5}{3}{3.5}{3}{0}
\fxrlink{2}{3}{2.5}{4}{2.5}{3}{0}
\fxrlink{2}{3}{3.5}{4}{3.5}{2}{0}
\fxrlink{2}{0}{0.25}{0}{0.75}{2}{0}
\fxrlink{2}{0}{0.75}{0}{1.25}{3}{0}
\fxrlink{2}{0}{1.25}{0}{1.75}{3}{0}
\fxrlink{2}{0}{1.75}{0}{2.5}{2}{0}
\fxrlink{2}{0}{2.5}{0}{3.5}{2}{0}
\fxrlink{2}{4}{0.5}{4}{1.5}{2}{0}
\fxrlink{2}{4}{1.5}{4}{2.5}{3}{0}
\fxrlink{2}{4}{2.5}{4}{3.5}{3}{0}
\fxrdof{2}{0.5}{0.25}{0}
\fxrdof{2}{0.5}{0.75}{0}
\fxrdof{2}{0.5}{1.25}{0}
\fxrdof{2}{0.5}{1.75}{0}
\fxrdof{2}{1}{0.25}{0}
\fxrdof{2}{1}{0.75}{0}
\fxrdof{2}{1}{1.25}{0}
\fxrdof{2}{1}{1.75}{0}
\fxrdof{2}{1}{2.5}{0}
\fxrdof{2}{1}{3.5}{0}
\fxrdof{2}{1.5}{0.25}{0}
\fxrdof{2}{1.5}{0.75}{0}
\fxrdof{2}{1.5}{1.25}{0}
\fxrdof{2}{1.5}{1.75}{0}
\fxrdof{2}{2}{0.25}{0}
\fxrdof{2}{2}{0.75}{0}
\fxrdof{2}{2}{1.25}{0}
\fxrdof{2}{2}{1.75}{0}
\fxrdof{2}{2}{2.5}{0}
\fxrdof{2}{2}{3.5}{0}
\fxrdof{2}{3}{0.5}{0}
\fxrdof{2}{3}{1.5}{0}
\fxrdof{2}{3}{2.5}{0}
\fxrdof{2}{3}{3.5}{0}
\fxrdof{2}{0}{0.25}{1}
\fxrdof{2}{0}{0.75}{1}
\fxrdof{2}{0}{1.25}{1}
\fxrdof{2}{0}{1.75}{1}
\fxrdof{2}{0}{2.5}{1}
\fxrdof{2}{0}{3.5}{1}
\fxrdof{2}{4}{0.5}{1}
\fxrdof{2}{4}{1.5}{1}
\fxrdof{2}{4}{2.5}{1}
\fxrdof{2}{4}{3.5}{1}
%
      \fxrnone
      \def\fxrcell##1##2##3##4##5{\draw[tr cell] (##1,##2) rectangle (##3,##4);}%
      \ifnum#1=2 \def\fxrplic##1##2##3##4{\draw[tr plic] (##1,##2) -- (##3,##4);}\fi
      % Generated by docs/techref/figures/generate.jl; do not edit.
\fxrcell{2}{0}{3}{1}{0}
\fxrliquid{(2,0) -- (3,0) -- (3,1) -- (2,1) -- cycle}
\fxrcell{3}{0}{4}{1}{0}
\fxrliquid{(3,0) -- (4,0) -- (4,1) -- (3,1) -- cycle}
\fxrcell{2}{1}{3}{2}{0}
\fxrliquid{(2,1) -- (3,1) -- (3,2) -- (2.571429,2) -- (2,1.8) -- cycle}
\fxrplic{2.571429}{2}{2}{1.8}
\fxrcell{3}{1}{4}{2}{0}
\fxrliquid{(3,1) -- (4,1) -- (4,2) -- (3,2) -- cycle}
\fxrcell{0}{2}{1}{3}{0}
\fxrcell{1}{2}{2}{3}{0}
\fxrcell{0}{3}{1}{4}{0}
\fxrcell{1}{3}{2}{4}{0}
\fxrcell{2}{2}{3}{3}{0}
\fxrliquid{(2.571429,2) -- (3,2) -- (3,2.15) -- cycle}
\fxrplic{2.571429}{2}{3}{2.15}
\fxrcell{3}{2}{4}{3}{0}
\fxrliquid{(3,2) -- (4,2) -- (4,2.5) -- (3,2.15) -- cycle}
\fxrplic{4}{2.5}{3}{2.15}
\fxrcell{2}{3}{3}{4}{0}
\fxrcell{3}{3}{4}{4}{0}
\fxrcell{0}{0}{0.5}{0.5}{1}
\fxrliquid{(0,0) -- (0.5,0) -- (0.5,0.5) -- (0,0.5) -- cycle}
\fxrcell{0.5}{0}{1}{0.5}{1}
\fxrliquid{(0.5,0) -- (1,0) -- (1,0.5) -- (0.5,0.5) -- cycle}
\fxrcell{0}{0.5}{0.5}{1}{1}
\fxrliquid{(0,0.5) -- (0.5,0.5) -- (0.5,1) -- (0,1) -- cycle}
\fxrcell{0.5}{0.5}{1}{1}{1}
\fxrliquid{(0.5,0.5) -- (1,0.5) -- (1,1) -- (0.5,1) -- cycle}
\fxrcell{1}{0}{1.5}{0.5}{1}
\fxrliquid{(1,0) -- (1.5,0) -- (1.5,0.5) -- (1,0.5) -- cycle}
\fxrcell{1.5}{0}{2}{0.5}{1}
\fxrliquid{(1.5,0) -- (2,0) -- (2,0.5) -- (1.5,0.5) -- cycle}
\fxrcell{1}{0.5}{1.5}{1}{1}
\fxrliquid{(1,0.5) -- (1.5,0.5) -- (1.5,1) -- (1,1) -- cycle}
\fxrcell{1.5}{0.5}{2}{1}{1}
\fxrliquid{(1.5,0.5) -- (2,0.5) -- (2,1) -- (1.5,1) -- cycle}
\fxrcell{0}{1}{0.5}{1.5}{1}
\fxrliquid{(0,1) -- (0.5,1) -- (0.5,1.275) -- (0,1.1) -- cycle}
\fxrplic{0.5}{1.275}{0}{1.1}
\fxrcell{0.5}{1}{1}{1.5}{1}
\fxrliquid{(0.5,1) -- (1,1) -- (1,1.45) -- (0.5,1.275) -- cycle}
\fxrplic{1}{1.45}{0.5}{1.275}
\fxrcell{0}{1.5}{0.5}{2}{1}
\fxrcell{0.5}{1.5}{1}{2}{1}
\fxrcell{1}{1}{1.5}{1.5}{1}
\fxrliquid{(1,1) -- (1.5,1) -- (1.5,1.5) -- (1.142857,1.5) -- (1,1.45) -- cycle}
\fxrplic{1.142857}{1.5}{1}{1.45}
\fxrcell{1.5}{1}{2}{1.5}{1}
\fxrliquid{(1.5,1) -- (2,1) -- (2,1.5) -- (1.5,1.5) -- cycle}
\fxrcell{1}{1.5}{1.5}{2}{1}
\fxrliquid{(1.142857,1.5) -- (1.5,1.5) -- (1.5,1.625) -- cycle}
\fxrplic{1.142857}{1.5}{1.5}{1.625}
\fxrcell{1.5}{1.5}{2}{2}{1}
\fxrliquid{(1.5,1.5) -- (2,1.5) -- (2,1.8) -- (1.5,1.625) -- cycle}
\fxrplic{2}{1.8}{1.5}{1.625}
\fxrlink{1}{0.5}{0.25}{0.5}{0.75}{0}{0}
\fxrlink{1}{0.5}{0.25}{1}{0.25}{0}{0}
\fxrlink{1}{0.5}{0.75}{0.5}{1.25}{0}{0}
\fxrlink{1}{0.5}{0.75}{1}{0.75}{0}{0}
\fxrlink{1}{0.5}{1.25}{0.5}{1.75}{0}{0}
\fxrlink{1}{0.5}{1.25}{1}{1.25}{0}{0}
\fxrlink{1}{0.5}{1.75}{1}{1.75}{0}{0}
\fxrlink{1}{0.5}{1.75}{1}{2.5}{1}{0}
\fxrlink{1}{1}{0.25}{1}{0.75}{0}{0}
\fxrlink{1}{1}{0.25}{1.5}{0.25}{0}{0}
\fxrlink{1}{1}{0.75}{1}{1.25}{0}{0}
\fxrlink{1}{1}{0.75}{1.5}{0.75}{0}{0}
\fxrlink{1}{1}{1.25}{1}{1.75}{0}{0}
\fxrlink{1}{1}{1.25}{1.5}{1.25}{0}{0}
\fxrlink{1}{1}{1.75}{1}{2.5}{1}{0}
\fxrlink{1}{1}{1.75}{1.5}{1.75}{0}{0}
\fxrlink{1}{1}{2.5}{1}{3.5}{0}{0}
\fxrlink{1}{1}{2.5}{1.5}{1.75}{1}{0}
\fxrlink{1}{1}{2.5}{2}{2.5}{0}{0}
\fxrlink{1}{1}{3.5}{2}{3.5}{0}{0}
\fxrlink{1}{1.5}{0.25}{1.5}{0.75}{0}{0}
\fxrlink{1}{1.5}{0.25}{2}{0.25}{1}{0}
\fxrlink{1}{1.5}{0.75}{1.5}{1.25}{0}{0}
\fxrlink{1}{1.5}{0.75}{2}{0.75}{1}{0}
\fxrlink{1}{1.5}{1.25}{1.5}{1.75}{0}{0}
\fxrlink{1}{1.5}{1.25}{2}{1.25}{1}{0}
\fxrlink{1}{1.5}{1.75}{2}{1.75}{1}{0}
\fxrlink{1}{1.5}{1.75}{2}{2.5}{1}{0}
\fxrlink{1}{2}{0.25}{2}{0.75}{1}{1}
\fxrlink{1}{2}{0.25}{2}{0.75}{1}{-1}
\fxrlink{1}{2}{0.25}{2}{1.25}{1}{0}
\fxrlink{1}{2}{0.25}{2}{1.75}{1}{0}
\fxrlink{1}{2}{0.25}{3}{0.5}{1}{0}
\fxrlink{1}{2}{0.75}{2}{1.25}{1}{0}
\fxrlink{1}{2}{0.75}{2}{1.75}{1}{0}
\fxrlink{1}{2}{0.75}{3}{0.5}{1}{0}
\fxrlink{1}{2}{1.25}{2}{1.75}{1}{1}
\fxrlink{1}{2}{1.25}{2}{1.75}{1}{-1}
\fxrlink{1}{2}{1.25}{2}{2.5}{1}{0}
\fxrlink{1}{2}{1.25}{3}{1.5}{1}{0}
\fxrlink{1}{2}{1.75}{2}{2.5}{1}{0}
\fxrlink{1}{2}{1.75}{3}{1.5}{1}{0}
\fxrlink{1}{2}{2.5}{2}{3.5}{0}{0}
\fxrlink{1}{2}{2.5}{3}{2.5}{0}{0}
\fxrlink{1}{2}{3.5}{3}{3.5}{0}{0}
\fxrlink{1}{3}{0.5}{3}{1.5}{0}{0}
\fxrlink{1}{3}{1.5}{3}{2.5}{0}{0}
\fxrlink{1}{3}{2.5}{3}{3.5}{0}{0}
\fxrlink{1}{0.5}{0.25}{0}{0.25}{4}{0}
\fxrlink{1}{0.5}{0.75}{0}{0.75}{4}{0}
\fxrlink{1}{0.5}{1.25}{0}{1.25}{4}{0}
\fxrlink{1}{0.5}{1.75}{0}{1.75}{4}{0}
\fxrlink{1}{1}{2.5}{0}{2.5}{4}{0}
\fxrlink{1}{1}{3.5}{0}{3.5}{4}{0}
\fxrlink{1}{3}{0.5}{4}{0.5}{4}{0}
\fxrlink{1}{3}{1.5}{4}{1.5}{4}{0}
\fxrlink{1}{3}{2.5}{4}{2.5}{4}{0}
\fxrlink{1}{3}{3.5}{4}{3.5}{4}{0}
\fxrlink{1}{0.5}{0.25}{0.25}{0}{5}{0}
\fxrlink{1}{0.5}{0.25}{0.75}{0}{5}{0}
\fxrlink{1}{1}{0.25}{0.75}{0}{5}{0}
\fxrlink{1}{1}{0.25}{1.25}{0}{5}{0}
\fxrlink{1}{1.5}{0.25}{1.25}{0}{5}{0}
\fxrlink{1}{1.5}{0.25}{1.75}{0}{5}{0}
\fxrlink{1}{2}{0.25}{1.75}{0}{5}{0}
\fxrlink{1}{2}{0.25}{2.5}{0}{5}{0}
\fxrlink{1}{2}{0.75}{2.5}{0}{5}{0}
\fxrlink{1}{3}{0.5}{2.5}{0}{5}{0}
\fxrlink{1}{3}{0.5}{3.5}{0}{5}{0}
\fxrlink{1}{1}{3.5}{0.5}{4}{5}{0}
\fxrlink{1}{1}{3.5}{1.5}{4}{5}{0}
\fxrlink{1}{2}{3.5}{1.5}{4}{5}{0}
\fxrlink{1}{2}{3.5}{2.5}{4}{5}{0}
\fxrlink{1}{3}{3.5}{2.5}{4}{5}{0}
\fxrlink{1}{3}{3.5}{3.5}{4}{5}{0}
\fxrlink{1}{0.5}{1.75}{0}{2.5}{6}{0}
\fxrdof{1}{0.5}{0.25}{0}
\fxrdof{1}{0.5}{0.75}{0}
\fxrdof{1}{0.5}{1.25}{0}
\fxrdof{1}{0.5}{1.75}{0}
\fxrdof{1}{1}{0.25}{0}
\fxrdof{1}{1}{0.75}{0}
\fxrdof{1}{1}{1.25}{0}
\fxrdof{1}{1}{1.75}{0}
\fxrdof{1}{1}{2.5}{0}
\fxrdof{1}{1}{3.5}{0}
\fxrdof{1}{1.5}{0.25}{0}
\fxrdof{1}{1.5}{0.75}{0}
\fxrdof{1}{1.5}{1.25}{0}
\fxrdof{1}{1.5}{1.75}{0}
\fxrdof{1}{2}{0.25}{0}
\fxrdof{1}{2}{0.75}{0}
\fxrdof{1}{2}{1.25}{0}
\fxrdof{1}{2}{1.75}{0}
\fxrdof{1}{2}{2.5}{0}
\fxrdof{1}{2}{3.5}{0}
\fxrdof{1}{3}{0.5}{0}
\fxrdof{1}{3}{1.5}{0}
\fxrdof{1}{3}{2.5}{0}
\fxrdof{1}{3}{3.5}{0}
\fxrlink{2}{0.5}{0.25}{0.5}{0.75}{0}{0}
\fxrlink{2}{0.5}{0.25}{1}{0.25}{0}{0}
\fxrlink{2}{0.5}{0.25}{0}{0.25}{2}{0}
\fxrlink{2}{0.5}{0.75}{0.5}{1.25}{3}{0}
\fxrlink{2}{0.5}{0.75}{1}{0.75}{0}{0}
\fxrlink{2}{0.5}{0.75}{0}{0.75}{2}{0}
\fxrlink{2}{0.5}{1.25}{0.5}{1.75}{3}{0}
\fxrlink{2}{0.5}{1.25}{1}{1.25}{3}{0}
\fxrlink{2}{0.5}{1.25}{0}{1.25}{3}{0}
\fxrlink{2}{0.5}{1.75}{1}{1.75}{3}{0}
\fxrlink{2}{0.5}{1.75}{1}{2.5}{1}{0}
\fxrlink{2}{0.5}{1.75}{0}{1.75}{2}{0}
\fxrlink{2}{0.5}{1.75}{0}{2.5}{2}{0}
\fxrlink{2}{1}{0.25}{1}{0.75}{0}{0}
\fxrlink{2}{1}{0.25}{1.5}{0.25}{0}{0}
\fxrlink{2}{1}{0.75}{1}{1.25}{3}{0}
\fxrlink{2}{1}{0.75}{1.5}{0.75}{0}{0}
\fxrlink{2}{1}{1.25}{1}{1.75}{3}{0}
\fxrlink{2}{1}{1.25}{1.5}{1.25}{3}{0}
\fxrlink{2}{1}{1.75}{1}{2.5}{3}{0}
\fxrlink{2}{1}{1.75}{1.5}{1.75}{3}{0}
\fxrlink{2}{1}{2.5}{1}{3.5}{0}{0}
\fxrlink{2}{1}{2.5}{1.5}{1.75}{3}{0}
\fxrlink{2}{1}{2.5}{2}{2.5}{3}{0}
\fxrlink{2}{1}{2.5}{0}{2.5}{2}{0}
\fxrlink{2}{1}{3.5}{2}{3.5}{0}{0}
\fxrlink{2}{1}{3.5}{0}{3.5}{2}{0}
\fxrlink{2}{1.5}{0.25}{1.5}{0.75}{0}{0}
\fxrlink{2}{1.5}{0.25}{2}{0.25}{1}{0}
\fxrlink{2}{1.5}{0.75}{1.5}{1.25}{3}{0}
\fxrlink{2}{1.5}{0.75}{2}{0.75}{1}{0}
\fxrlink{2}{1.5}{1.25}{1.5}{1.75}{3}{0}
\fxrlink{2}{1.5}{1.25}{2}{1.25}{3}{0}
\fxrlink{2}{1.5}{1.75}{2}{1.75}{3}{0}
\fxrlink{2}{1.5}{1.75}{2}{2.5}{3}{0}
\fxrlink{2}{2}{0.25}{2}{0.75}{1}{1}
\fxrlink{2}{2}{0.25}{2}{0.75}{1}{-1}
\fxrlink{2}{2}{0.25}{2}{1.25}{3}{0}
\fxrlink{2}{2}{0.25}{2}{1.75}{3}{0}
\fxrlink{2}{2}{0.25}{3}{0.5}{1}{0}
\fxrlink{2}{2}{0.75}{2}{1.25}{3}{0}
\fxrlink{2}{2}{0.75}{2}{1.75}{3}{0}
\fxrlink{2}{2}{0.75}{3}{0.5}{1}{0}
\fxrlink{2}{2}{1.25}{2}{1.75}{3}{1}
\fxrlink{2}{2}{1.25}{2}{1.75}{3}{-1}
\fxrlink{2}{2}{1.25}{2}{2.5}{3}{0}
\fxrlink{2}{2}{1.25}{3}{1.5}{3}{0}
\fxrlink{2}{2}{1.75}{2}{2.5}{3}{0}
\fxrlink{2}{2}{1.75}{3}{1.5}{3}{0}
\fxrlink{2}{2}{2.5}{2}{3.5}{3}{0}
\fxrlink{2}{2}{2.5}{3}{2.5}{3}{0}
\fxrlink{2}{2}{3.5}{3}{3.5}{0}{0}
\fxrlink{2}{3}{0.5}{3}{1.5}{3}{0}
\fxrlink{2}{3}{0.5}{4}{0.5}{2}{0}
\fxrlink{2}{3}{1.5}{3}{2.5}{3}{0}
\fxrlink{2}{3}{1.5}{4}{1.5}{3}{0}
\fxrlink{2}{3}{2.5}{3}{3.5}{3}{0}
\fxrlink{2}{3}{2.5}{4}{2.5}{3}{0}
\fxrlink{2}{3}{3.5}{4}{3.5}{2}{0}
\fxrlink{2}{0}{0.25}{0}{0.75}{2}{0}
\fxrlink{2}{0}{0.75}{0}{1.25}{3}{0}
\fxrlink{2}{0}{1.25}{0}{1.75}{3}{0}
\fxrlink{2}{0}{1.75}{0}{2.5}{2}{0}
\fxrlink{2}{0}{2.5}{0}{3.5}{2}{0}
\fxrlink{2}{4}{0.5}{4}{1.5}{2}{0}
\fxrlink{2}{4}{1.5}{4}{2.5}{3}{0}
\fxrlink{2}{4}{2.5}{4}{3.5}{3}{0}
\fxrdof{2}{0.5}{0.25}{0}
\fxrdof{2}{0.5}{0.75}{0}
\fxrdof{2}{0.5}{1.25}{0}
\fxrdof{2}{0.5}{1.75}{0}
\fxrdof{2}{1}{0.25}{0}
\fxrdof{2}{1}{0.75}{0}
\fxrdof{2}{1}{1.25}{0}
\fxrdof{2}{1}{1.75}{0}
\fxrdof{2}{1}{2.5}{0}
\fxrdof{2}{1}{3.5}{0}
\fxrdof{2}{1.5}{0.25}{0}
\fxrdof{2}{1.5}{0.75}{0}
\fxrdof{2}{1.5}{1.25}{0}
\fxrdof{2}{1.5}{1.75}{0}
\fxrdof{2}{2}{0.25}{0}
\fxrdof{2}{2}{0.75}{0}
\fxrdof{2}{2}{1.25}{0}
\fxrdof{2}{2}{1.75}{0}
\fxrdof{2}{2}{2.5}{0}
\fxrdof{2}{2}{3.5}{0}
\fxrdof{2}{3}{0.5}{0}
\fxrdof{2}{3}{1.5}{0}
\fxrdof{2}{3}{2.5}{0}
\fxrdof{2}{3}{3.5}{0}
\fxrdof{2}{0}{0.25}{1}
\fxrdof{2}{0}{0.75}{1}
\fxrdof{2}{0}{1.25}{1}
\fxrdof{2}{0}{1.75}{1}
\fxrdof{2}{0}{2.5}{1}
\fxrdof{2}{0}{3.5}{1}
\fxrdof{2}{4}{0.5}{1}
\fxrdof{2}{4}{1.5}{1}
\fxrdof{2}{4}{2.5}{1}
\fxrdof{2}{4}{3.5}{1}
%
      \fill[tr wall] (-0.14,-0.14) rectangle (4.14,0);
      \fill[tr wall] (-0.14,4) rectangle (4.14,4.14);
      \fill[tr wall] (-0.14,0) rectangle (0,4);
      \fill[tr wall] (4,0) rectangle (4.14,4);
      \draw[black, line width=0.7pt] (0,0) rectangle (4,4);
      \fxrnone
      \def\fxrlink##1##2##3##4##5##6##7{\ifnum##1=#1
        \draw[fxr ##6] (##2,##3) to[fxr bend ##7] (##4,##5);\fi}%
      % Generated by docs/techref/figures/generate.jl; do not edit.
\fxrcell{2}{0}{3}{1}{0}
\fxrliquid{(2,0) -- (3,0) -- (3,1) -- (2,1) -- cycle}
\fxrcell{3}{0}{4}{1}{0}
\fxrliquid{(3,0) -- (4,0) -- (4,1) -- (3,1) -- cycle}
\fxrcell{2}{1}{3}{2}{0}
\fxrliquid{(2,1) -- (3,1) -- (3,2) -- (2.571429,2) -- (2,1.8) -- cycle}
\fxrplic{2.571429}{2}{2}{1.8}
\fxrcell{3}{1}{4}{2}{0}
\fxrliquid{(3,1) -- (4,1) -- (4,2) -- (3,2) -- cycle}
\fxrcell{0}{2}{1}{3}{0}
\fxrcell{1}{2}{2}{3}{0}
\fxrcell{0}{3}{1}{4}{0}
\fxrcell{1}{3}{2}{4}{0}
\fxrcell{2}{2}{3}{3}{0}
\fxrliquid{(2.571429,2) -- (3,2) -- (3,2.15) -- cycle}
\fxrplic{2.571429}{2}{3}{2.15}
\fxrcell{3}{2}{4}{3}{0}
\fxrliquid{(3,2) -- (4,2) -- (4,2.5) -- (3,2.15) -- cycle}
\fxrplic{4}{2.5}{3}{2.15}
\fxrcell{2}{3}{3}{4}{0}
\fxrcell{3}{3}{4}{4}{0}
\fxrcell{0}{0}{0.5}{0.5}{1}
\fxrliquid{(0,0) -- (0.5,0) -- (0.5,0.5) -- (0,0.5) -- cycle}
\fxrcell{0.5}{0}{1}{0.5}{1}
\fxrliquid{(0.5,0) -- (1,0) -- (1,0.5) -- (0.5,0.5) -- cycle}
\fxrcell{0}{0.5}{0.5}{1}{1}
\fxrliquid{(0,0.5) -- (0.5,0.5) -- (0.5,1) -- (0,1) -- cycle}
\fxrcell{0.5}{0.5}{1}{1}{1}
\fxrliquid{(0.5,0.5) -- (1,0.5) -- (1,1) -- (0.5,1) -- cycle}
\fxrcell{1}{0}{1.5}{0.5}{1}
\fxrliquid{(1,0) -- (1.5,0) -- (1.5,0.5) -- (1,0.5) -- cycle}
\fxrcell{1.5}{0}{2}{0.5}{1}
\fxrliquid{(1.5,0) -- (2,0) -- (2,0.5) -- (1.5,0.5) -- cycle}
\fxrcell{1}{0.5}{1.5}{1}{1}
\fxrliquid{(1,0.5) -- (1.5,0.5) -- (1.5,1) -- (1,1) -- cycle}
\fxrcell{1.5}{0.5}{2}{1}{1}
\fxrliquid{(1.5,0.5) -- (2,0.5) -- (2,1) -- (1.5,1) -- cycle}
\fxrcell{0}{1}{0.5}{1.5}{1}
\fxrliquid{(0,1) -- (0.5,1) -- (0.5,1.275) -- (0,1.1) -- cycle}
\fxrplic{0.5}{1.275}{0}{1.1}
\fxrcell{0.5}{1}{1}{1.5}{1}
\fxrliquid{(0.5,1) -- (1,1) -- (1,1.45) -- (0.5,1.275) -- cycle}
\fxrplic{1}{1.45}{0.5}{1.275}
\fxrcell{0}{1.5}{0.5}{2}{1}
\fxrcell{0.5}{1.5}{1}{2}{1}
\fxrcell{1}{1}{1.5}{1.5}{1}
\fxrliquid{(1,1) -- (1.5,1) -- (1.5,1.5) -- (1.142857,1.5) -- (1,1.45) -- cycle}
\fxrplic{1.142857}{1.5}{1}{1.45}
\fxrcell{1.5}{1}{2}{1.5}{1}
\fxrliquid{(1.5,1) -- (2,1) -- (2,1.5) -- (1.5,1.5) -- cycle}
\fxrcell{1}{1.5}{1.5}{2}{1}
\fxrliquid{(1.142857,1.5) -- (1.5,1.5) -- (1.5,1.625) -- cycle}
\fxrplic{1.142857}{1.5}{1.5}{1.625}
\fxrcell{1.5}{1.5}{2}{2}{1}
\fxrliquid{(1.5,1.5) -- (2,1.5) -- (2,1.8) -- (1.5,1.625) -- cycle}
\fxrplic{2}{1.8}{1.5}{1.625}
\fxrlink{1}{0.5}{0.25}{0.5}{0.75}{0}{0}
\fxrlink{1}{0.5}{0.25}{1}{0.25}{0}{0}
\fxrlink{1}{0.5}{0.75}{0.5}{1.25}{0}{0}
\fxrlink{1}{0.5}{0.75}{1}{0.75}{0}{0}
\fxrlink{1}{0.5}{1.25}{0.5}{1.75}{0}{0}
\fxrlink{1}{0.5}{1.25}{1}{1.25}{0}{0}
\fxrlink{1}{0.5}{1.75}{1}{1.75}{0}{0}
\fxrlink{1}{0.5}{1.75}{1}{2.5}{1}{0}
\fxrlink{1}{1}{0.25}{1}{0.75}{0}{0}
\fxrlink{1}{1}{0.25}{1.5}{0.25}{0}{0}
\fxrlink{1}{1}{0.75}{1}{1.25}{0}{0}
\fxrlink{1}{1}{0.75}{1.5}{0.75}{0}{0}
\fxrlink{1}{1}{1.25}{1}{1.75}{0}{0}
\fxrlink{1}{1}{1.25}{1.5}{1.25}{0}{0}
\fxrlink{1}{1}{1.75}{1}{2.5}{1}{0}
\fxrlink{1}{1}{1.75}{1.5}{1.75}{0}{0}
\fxrlink{1}{1}{2.5}{1}{3.5}{0}{0}
\fxrlink{1}{1}{2.5}{1.5}{1.75}{1}{0}
\fxrlink{1}{1}{2.5}{2}{2.5}{0}{0}
\fxrlink{1}{1}{3.5}{2}{3.5}{0}{0}
\fxrlink{1}{1.5}{0.25}{1.5}{0.75}{0}{0}
\fxrlink{1}{1.5}{0.25}{2}{0.25}{1}{0}
\fxrlink{1}{1.5}{0.75}{1.5}{1.25}{0}{0}
\fxrlink{1}{1.5}{0.75}{2}{0.75}{1}{0}
\fxrlink{1}{1.5}{1.25}{1.5}{1.75}{0}{0}
\fxrlink{1}{1.5}{1.25}{2}{1.25}{1}{0}
\fxrlink{1}{1.5}{1.75}{2}{1.75}{1}{0}
\fxrlink{1}{1.5}{1.75}{2}{2.5}{1}{0}
\fxrlink{1}{2}{0.25}{2}{0.75}{1}{1}
\fxrlink{1}{2}{0.25}{2}{0.75}{1}{-1}
\fxrlink{1}{2}{0.25}{2}{1.25}{1}{0}
\fxrlink{1}{2}{0.25}{2}{1.75}{1}{0}
\fxrlink{1}{2}{0.25}{3}{0.5}{1}{0}
\fxrlink{1}{2}{0.75}{2}{1.25}{1}{0}
\fxrlink{1}{2}{0.75}{2}{1.75}{1}{0}
\fxrlink{1}{2}{0.75}{3}{0.5}{1}{0}
\fxrlink{1}{2}{1.25}{2}{1.75}{1}{1}
\fxrlink{1}{2}{1.25}{2}{1.75}{1}{-1}
\fxrlink{1}{2}{1.25}{2}{2.5}{1}{0}
\fxrlink{1}{2}{1.25}{3}{1.5}{1}{0}
\fxrlink{1}{2}{1.75}{2}{2.5}{1}{0}
\fxrlink{1}{2}{1.75}{3}{1.5}{1}{0}
\fxrlink{1}{2}{2.5}{2}{3.5}{0}{0}
\fxrlink{1}{2}{2.5}{3}{2.5}{0}{0}
\fxrlink{1}{2}{3.5}{3}{3.5}{0}{0}
\fxrlink{1}{3}{0.5}{3}{1.5}{0}{0}
\fxrlink{1}{3}{1.5}{3}{2.5}{0}{0}
\fxrlink{1}{3}{2.5}{3}{3.5}{0}{0}
\fxrlink{1}{0.5}{0.25}{0}{0.25}{4}{0}
\fxrlink{1}{0.5}{0.75}{0}{0.75}{4}{0}
\fxrlink{1}{0.5}{1.25}{0}{1.25}{4}{0}
\fxrlink{1}{0.5}{1.75}{0}{1.75}{4}{0}
\fxrlink{1}{1}{2.5}{0}{2.5}{4}{0}
\fxrlink{1}{1}{3.5}{0}{3.5}{4}{0}
\fxrlink{1}{3}{0.5}{4}{0.5}{4}{0}
\fxrlink{1}{3}{1.5}{4}{1.5}{4}{0}
\fxrlink{1}{3}{2.5}{4}{2.5}{4}{0}
\fxrlink{1}{3}{3.5}{4}{3.5}{4}{0}
\fxrlink{1}{0.5}{0.25}{0.25}{0}{5}{0}
\fxrlink{1}{0.5}{0.25}{0.75}{0}{5}{0}
\fxrlink{1}{1}{0.25}{0.75}{0}{5}{0}
\fxrlink{1}{1}{0.25}{1.25}{0}{5}{0}
\fxrlink{1}{1.5}{0.25}{1.25}{0}{5}{0}
\fxrlink{1}{1.5}{0.25}{1.75}{0}{5}{0}
\fxrlink{1}{2}{0.25}{1.75}{0}{5}{0}
\fxrlink{1}{2}{0.25}{2.5}{0}{5}{0}
\fxrlink{1}{2}{0.75}{2.5}{0}{5}{0}
\fxrlink{1}{3}{0.5}{2.5}{0}{5}{0}
\fxrlink{1}{3}{0.5}{3.5}{0}{5}{0}
\fxrlink{1}{1}{3.5}{0.5}{4}{5}{0}
\fxrlink{1}{1}{3.5}{1.5}{4}{5}{0}
\fxrlink{1}{2}{3.5}{1.5}{4}{5}{0}
\fxrlink{1}{2}{3.5}{2.5}{4}{5}{0}
\fxrlink{1}{3}{3.5}{2.5}{4}{5}{0}
\fxrlink{1}{3}{3.5}{3.5}{4}{5}{0}
\fxrlink{1}{0.5}{1.75}{0}{2.5}{6}{0}
\fxrdof{1}{0.5}{0.25}{0}
\fxrdof{1}{0.5}{0.75}{0}
\fxrdof{1}{0.5}{1.25}{0}
\fxrdof{1}{0.5}{1.75}{0}
\fxrdof{1}{1}{0.25}{0}
\fxrdof{1}{1}{0.75}{0}
\fxrdof{1}{1}{1.25}{0}
\fxrdof{1}{1}{1.75}{0}
\fxrdof{1}{1}{2.5}{0}
\fxrdof{1}{1}{3.5}{0}
\fxrdof{1}{1.5}{0.25}{0}
\fxrdof{1}{1.5}{0.75}{0}
\fxrdof{1}{1.5}{1.25}{0}
\fxrdof{1}{1.5}{1.75}{0}
\fxrdof{1}{2}{0.25}{0}
\fxrdof{1}{2}{0.75}{0}
\fxrdof{1}{2}{1.25}{0}
\fxrdof{1}{2}{1.75}{0}
\fxrdof{1}{2}{2.5}{0}
\fxrdof{1}{2}{3.5}{0}
\fxrdof{1}{3}{0.5}{0}
\fxrdof{1}{3}{1.5}{0}
\fxrdof{1}{3}{2.5}{0}
\fxrdof{1}{3}{3.5}{0}
\fxrlink{2}{0.5}{0.25}{0.5}{0.75}{0}{0}
\fxrlink{2}{0.5}{0.25}{1}{0.25}{0}{0}
\fxrlink{2}{0.5}{0.25}{0}{0.25}{2}{0}
\fxrlink{2}{0.5}{0.75}{0.5}{1.25}{3}{0}
\fxrlink{2}{0.5}{0.75}{1}{0.75}{0}{0}
\fxrlink{2}{0.5}{0.75}{0}{0.75}{2}{0}
\fxrlink{2}{0.5}{1.25}{0.5}{1.75}{3}{0}
\fxrlink{2}{0.5}{1.25}{1}{1.25}{3}{0}
\fxrlink{2}{0.5}{1.25}{0}{1.25}{3}{0}
\fxrlink{2}{0.5}{1.75}{1}{1.75}{3}{0}
\fxrlink{2}{0.5}{1.75}{1}{2.5}{1}{0}
\fxrlink{2}{0.5}{1.75}{0}{1.75}{2}{0}
\fxrlink{2}{0.5}{1.75}{0}{2.5}{2}{0}
\fxrlink{2}{1}{0.25}{1}{0.75}{0}{0}
\fxrlink{2}{1}{0.25}{1.5}{0.25}{0}{0}
\fxrlink{2}{1}{0.75}{1}{1.25}{3}{0}
\fxrlink{2}{1}{0.75}{1.5}{0.75}{0}{0}
\fxrlink{2}{1}{1.25}{1}{1.75}{3}{0}
\fxrlink{2}{1}{1.25}{1.5}{1.25}{3}{0}
\fxrlink{2}{1}{1.75}{1}{2.5}{3}{0}
\fxrlink{2}{1}{1.75}{1.5}{1.75}{3}{0}
\fxrlink{2}{1}{2.5}{1}{3.5}{0}{0}
\fxrlink{2}{1}{2.5}{1.5}{1.75}{3}{0}
\fxrlink{2}{1}{2.5}{2}{2.5}{3}{0}
\fxrlink{2}{1}{2.5}{0}{2.5}{2}{0}
\fxrlink{2}{1}{3.5}{2}{3.5}{0}{0}
\fxrlink{2}{1}{3.5}{0}{3.5}{2}{0}
\fxrlink{2}{1.5}{0.25}{1.5}{0.75}{0}{0}
\fxrlink{2}{1.5}{0.25}{2}{0.25}{1}{0}
\fxrlink{2}{1.5}{0.75}{1.5}{1.25}{3}{0}
\fxrlink{2}{1.5}{0.75}{2}{0.75}{1}{0}
\fxrlink{2}{1.5}{1.25}{1.5}{1.75}{3}{0}
\fxrlink{2}{1.5}{1.25}{2}{1.25}{3}{0}
\fxrlink{2}{1.5}{1.75}{2}{1.75}{3}{0}
\fxrlink{2}{1.5}{1.75}{2}{2.5}{3}{0}
\fxrlink{2}{2}{0.25}{2}{0.75}{1}{1}
\fxrlink{2}{2}{0.25}{2}{0.75}{1}{-1}
\fxrlink{2}{2}{0.25}{2}{1.25}{3}{0}
\fxrlink{2}{2}{0.25}{2}{1.75}{3}{0}
\fxrlink{2}{2}{0.25}{3}{0.5}{1}{0}
\fxrlink{2}{2}{0.75}{2}{1.25}{3}{0}
\fxrlink{2}{2}{0.75}{2}{1.75}{3}{0}
\fxrlink{2}{2}{0.75}{3}{0.5}{1}{0}
\fxrlink{2}{2}{1.25}{2}{1.75}{3}{1}
\fxrlink{2}{2}{1.25}{2}{1.75}{3}{-1}
\fxrlink{2}{2}{1.25}{2}{2.5}{3}{0}
\fxrlink{2}{2}{1.25}{3}{1.5}{3}{0}
\fxrlink{2}{2}{1.75}{2}{2.5}{3}{0}
\fxrlink{2}{2}{1.75}{3}{1.5}{3}{0}
\fxrlink{2}{2}{2.5}{2}{3.5}{3}{0}
\fxrlink{2}{2}{2.5}{3}{2.5}{3}{0}
\fxrlink{2}{2}{3.5}{3}{3.5}{0}{0}
\fxrlink{2}{3}{0.5}{3}{1.5}{3}{0}
\fxrlink{2}{3}{0.5}{4}{0.5}{2}{0}
\fxrlink{2}{3}{1.5}{3}{2.5}{3}{0}
\fxrlink{2}{3}{1.5}{4}{1.5}{3}{0}
\fxrlink{2}{3}{2.5}{3}{3.5}{3}{0}
\fxrlink{2}{3}{2.5}{4}{2.5}{3}{0}
\fxrlink{2}{3}{3.5}{4}{3.5}{2}{0}
\fxrlink{2}{0}{0.25}{0}{0.75}{2}{0}
\fxrlink{2}{0}{0.75}{0}{1.25}{3}{0}
\fxrlink{2}{0}{1.25}{0}{1.75}{3}{0}
\fxrlink{2}{0}{1.75}{0}{2.5}{2}{0}
\fxrlink{2}{0}{2.5}{0}{3.5}{2}{0}
\fxrlink{2}{4}{0.5}{4}{1.5}{2}{0}
\fxrlink{2}{4}{1.5}{4}{2.5}{3}{0}
\fxrlink{2}{4}{2.5}{4}{3.5}{3}{0}
\fxrdof{2}{0.5}{0.25}{0}
\fxrdof{2}{0.5}{0.75}{0}
\fxrdof{2}{0.5}{1.25}{0}
\fxrdof{2}{0.5}{1.75}{0}
\fxrdof{2}{1}{0.25}{0}
\fxrdof{2}{1}{0.75}{0}
\fxrdof{2}{1}{1.25}{0}
\fxrdof{2}{1}{1.75}{0}
\fxrdof{2}{1}{2.5}{0}
\fxrdof{2}{1}{3.5}{0}
\fxrdof{2}{1.5}{0.25}{0}
\fxrdof{2}{1.5}{0.75}{0}
\fxrdof{2}{1.5}{1.25}{0}
\fxrdof{2}{1.5}{1.75}{0}
\fxrdof{2}{2}{0.25}{0}
\fxrdof{2}{2}{0.75}{0}
\fxrdof{2}{2}{1.25}{0}
\fxrdof{2}{2}{1.75}{0}
\fxrdof{2}{2}{2.5}{0}
\fxrdof{2}{2}{3.5}{0}
\fxrdof{2}{3}{0.5}{0}
\fxrdof{2}{3}{1.5}{0}
\fxrdof{2}{3}{2.5}{0}
\fxrdof{2}{3}{3.5}{0}
\fxrdof{2}{0}{0.25}{1}
\fxrdof{2}{0}{0.75}{1}
\fxrdof{2}{0}{1.25}{1}
\fxrdof{2}{0}{1.75}{1}
\fxrdof{2}{0}{2.5}{1}
\fxrdof{2}{0}{3.5}{1}
\fxrdof{2}{4}{0.5}{1}
\fxrdof{2}{4}{1.5}{1}
\fxrdof{2}{4}{2.5}{1}
\fxrdof{2}{4}{3.5}{1}
%
      \fxrnone
      \def\fxrdof##1##2##3##4{\ifnum##1=#1
        \ifnum##4=1 \node[fxr bdof] at (##2,##3) {};\else \node[fxr dof] at (##2,##3) {};\fi\fi}%
      % Generated by docs/techref/figures/generate.jl; do not edit.
\fxrcell{2}{0}{3}{1}{0}
\fxrliquid{(2,0) -- (3,0) -- (3,1) -- (2,1) -- cycle}
\fxrcell{3}{0}{4}{1}{0}
\fxrliquid{(3,0) -- (4,0) -- (4,1) -- (3,1) -- cycle}
\fxrcell{2}{1}{3}{2}{0}
\fxrliquid{(2,1) -- (3,1) -- (3,2) -- (2.571429,2) -- (2,1.8) -- cycle}
\fxrplic{2.571429}{2}{2}{1.8}
\fxrcell{3}{1}{4}{2}{0}
\fxrliquid{(3,1) -- (4,1) -- (4,2) -- (3,2) -- cycle}
\fxrcell{0}{2}{1}{3}{0}
\fxrcell{1}{2}{2}{3}{0}
\fxrcell{0}{3}{1}{4}{0}
\fxrcell{1}{3}{2}{4}{0}
\fxrcell{2}{2}{3}{3}{0}
\fxrliquid{(2.571429,2) -- (3,2) -- (3,2.15) -- cycle}
\fxrplic{2.571429}{2}{3}{2.15}
\fxrcell{3}{2}{4}{3}{0}
\fxrliquid{(3,2) -- (4,2) -- (4,2.5) -- (3,2.15) -- cycle}
\fxrplic{4}{2.5}{3}{2.15}
\fxrcell{2}{3}{3}{4}{0}
\fxrcell{3}{3}{4}{4}{0}
\fxrcell{0}{0}{0.5}{0.5}{1}
\fxrliquid{(0,0) -- (0.5,0) -- (0.5,0.5) -- (0,0.5) -- cycle}
\fxrcell{0.5}{0}{1}{0.5}{1}
\fxrliquid{(0.5,0) -- (1,0) -- (1,0.5) -- (0.5,0.5) -- cycle}
\fxrcell{0}{0.5}{0.5}{1}{1}
\fxrliquid{(0,0.5) -- (0.5,0.5) -- (0.5,1) -- (0,1) -- cycle}
\fxrcell{0.5}{0.5}{1}{1}{1}
\fxrliquid{(0.5,0.5) -- (1,0.5) -- (1,1) -- (0.5,1) -- cycle}
\fxrcell{1}{0}{1.5}{0.5}{1}
\fxrliquid{(1,0) -- (1.5,0) -- (1.5,0.5) -- (1,0.5) -- cycle}
\fxrcell{1.5}{0}{2}{0.5}{1}
\fxrliquid{(1.5,0) -- (2,0) -- (2,0.5) -- (1.5,0.5) -- cycle}
\fxrcell{1}{0.5}{1.5}{1}{1}
\fxrliquid{(1,0.5) -- (1.5,0.5) -- (1.5,1) -- (1,1) -- cycle}
\fxrcell{1.5}{0.5}{2}{1}{1}
\fxrliquid{(1.5,0.5) -- (2,0.5) -- (2,1) -- (1.5,1) -- cycle}
\fxrcell{0}{1}{0.5}{1.5}{1}
\fxrliquid{(0,1) -- (0.5,1) -- (0.5,1.275) -- (0,1.1) -- cycle}
\fxrplic{0.5}{1.275}{0}{1.1}
\fxrcell{0.5}{1}{1}{1.5}{1}
\fxrliquid{(0.5,1) -- (1,1) -- (1,1.45) -- (0.5,1.275) -- cycle}
\fxrplic{1}{1.45}{0.5}{1.275}
\fxrcell{0}{1.5}{0.5}{2}{1}
\fxrcell{0.5}{1.5}{1}{2}{1}
\fxrcell{1}{1}{1.5}{1.5}{1}
\fxrliquid{(1,1) -- (1.5,1) -- (1.5,1.5) -- (1.142857,1.5) -- (1,1.45) -- cycle}
\fxrplic{1.142857}{1.5}{1}{1.45}
\fxrcell{1.5}{1}{2}{1.5}{1}
\fxrliquid{(1.5,1) -- (2,1) -- (2,1.5) -- (1.5,1.5) -- cycle}
\fxrcell{1}{1.5}{1.5}{2}{1}
\fxrliquid{(1.142857,1.5) -- (1.5,1.5) -- (1.5,1.625) -- cycle}
\fxrplic{1.142857}{1.5}{1.5}{1.625}
\fxrcell{1.5}{1.5}{2}{2}{1}
\fxrliquid{(1.5,1.5) -- (2,1.5) -- (2,1.8) -- (1.5,1.625) -- cycle}
\fxrplic{2}{1.8}{1.5}{1.625}
\fxrlink{1}{0.5}{0.25}{0.5}{0.75}{0}{0}
\fxrlink{1}{0.5}{0.25}{1}{0.25}{0}{0}
\fxrlink{1}{0.5}{0.75}{0.5}{1.25}{0}{0}
\fxrlink{1}{0.5}{0.75}{1}{0.75}{0}{0}
\fxrlink{1}{0.5}{1.25}{0.5}{1.75}{0}{0}
\fxrlink{1}{0.5}{1.25}{1}{1.25}{0}{0}
\fxrlink{1}{0.5}{1.75}{1}{1.75}{0}{0}
\fxrlink{1}{0.5}{1.75}{1}{2.5}{1}{0}
\fxrlink{1}{1}{0.25}{1}{0.75}{0}{0}
\fxrlink{1}{1}{0.25}{1.5}{0.25}{0}{0}
\fxrlink{1}{1}{0.75}{1}{1.25}{0}{0}
\fxrlink{1}{1}{0.75}{1.5}{0.75}{0}{0}
\fxrlink{1}{1}{1.25}{1}{1.75}{0}{0}
\fxrlink{1}{1}{1.25}{1.5}{1.25}{0}{0}
\fxrlink{1}{1}{1.75}{1}{2.5}{1}{0}
\fxrlink{1}{1}{1.75}{1.5}{1.75}{0}{0}
\fxrlink{1}{1}{2.5}{1}{3.5}{0}{0}
\fxrlink{1}{1}{2.5}{1.5}{1.75}{1}{0}
\fxrlink{1}{1}{2.5}{2}{2.5}{0}{0}
\fxrlink{1}{1}{3.5}{2}{3.5}{0}{0}
\fxrlink{1}{1.5}{0.25}{1.5}{0.75}{0}{0}
\fxrlink{1}{1.5}{0.25}{2}{0.25}{1}{0}
\fxrlink{1}{1.5}{0.75}{1.5}{1.25}{0}{0}
\fxrlink{1}{1.5}{0.75}{2}{0.75}{1}{0}
\fxrlink{1}{1.5}{1.25}{1.5}{1.75}{0}{0}
\fxrlink{1}{1.5}{1.25}{2}{1.25}{1}{0}
\fxrlink{1}{1.5}{1.75}{2}{1.75}{1}{0}
\fxrlink{1}{1.5}{1.75}{2}{2.5}{1}{0}
\fxrlink{1}{2}{0.25}{2}{0.75}{1}{1}
\fxrlink{1}{2}{0.25}{2}{0.75}{1}{-1}
\fxrlink{1}{2}{0.25}{2}{1.25}{1}{0}
\fxrlink{1}{2}{0.25}{2}{1.75}{1}{0}
\fxrlink{1}{2}{0.25}{3}{0.5}{1}{0}
\fxrlink{1}{2}{0.75}{2}{1.25}{1}{0}
\fxrlink{1}{2}{0.75}{2}{1.75}{1}{0}
\fxrlink{1}{2}{0.75}{3}{0.5}{1}{0}
\fxrlink{1}{2}{1.25}{2}{1.75}{1}{1}
\fxrlink{1}{2}{1.25}{2}{1.75}{1}{-1}
\fxrlink{1}{2}{1.25}{2}{2.5}{1}{0}
\fxrlink{1}{2}{1.25}{3}{1.5}{1}{0}
\fxrlink{1}{2}{1.75}{2}{2.5}{1}{0}
\fxrlink{1}{2}{1.75}{3}{1.5}{1}{0}
\fxrlink{1}{2}{2.5}{2}{3.5}{0}{0}
\fxrlink{1}{2}{2.5}{3}{2.5}{0}{0}
\fxrlink{1}{2}{3.5}{3}{3.5}{0}{0}
\fxrlink{1}{3}{0.5}{3}{1.5}{0}{0}
\fxrlink{1}{3}{1.5}{3}{2.5}{0}{0}
\fxrlink{1}{3}{2.5}{3}{3.5}{0}{0}
\fxrlink{1}{0.5}{0.25}{0}{0.25}{4}{0}
\fxrlink{1}{0.5}{0.75}{0}{0.75}{4}{0}
\fxrlink{1}{0.5}{1.25}{0}{1.25}{4}{0}
\fxrlink{1}{0.5}{1.75}{0}{1.75}{4}{0}
\fxrlink{1}{1}{2.5}{0}{2.5}{4}{0}
\fxrlink{1}{1}{3.5}{0}{3.5}{4}{0}
\fxrlink{1}{3}{0.5}{4}{0.5}{4}{0}
\fxrlink{1}{3}{1.5}{4}{1.5}{4}{0}
\fxrlink{1}{3}{2.5}{4}{2.5}{4}{0}
\fxrlink{1}{3}{3.5}{4}{3.5}{4}{0}
\fxrlink{1}{0.5}{0.25}{0.25}{0}{5}{0}
\fxrlink{1}{0.5}{0.25}{0.75}{0}{5}{0}
\fxrlink{1}{1}{0.25}{0.75}{0}{5}{0}
\fxrlink{1}{1}{0.25}{1.25}{0}{5}{0}
\fxrlink{1}{1.5}{0.25}{1.25}{0}{5}{0}
\fxrlink{1}{1.5}{0.25}{1.75}{0}{5}{0}
\fxrlink{1}{2}{0.25}{1.75}{0}{5}{0}
\fxrlink{1}{2}{0.25}{2.5}{0}{5}{0}
\fxrlink{1}{2}{0.75}{2.5}{0}{5}{0}
\fxrlink{1}{3}{0.5}{2.5}{0}{5}{0}
\fxrlink{1}{3}{0.5}{3.5}{0}{5}{0}
\fxrlink{1}{1}{3.5}{0.5}{4}{5}{0}
\fxrlink{1}{1}{3.5}{1.5}{4}{5}{0}
\fxrlink{1}{2}{3.5}{1.5}{4}{5}{0}
\fxrlink{1}{2}{3.5}{2.5}{4}{5}{0}
\fxrlink{1}{3}{3.5}{2.5}{4}{5}{0}
\fxrlink{1}{3}{3.5}{3.5}{4}{5}{0}
\fxrlink{1}{0.5}{1.75}{0}{2.5}{6}{0}
\fxrdof{1}{0.5}{0.25}{0}
\fxrdof{1}{0.5}{0.75}{0}
\fxrdof{1}{0.5}{1.25}{0}
\fxrdof{1}{0.5}{1.75}{0}
\fxrdof{1}{1}{0.25}{0}
\fxrdof{1}{1}{0.75}{0}
\fxrdof{1}{1}{1.25}{0}
\fxrdof{1}{1}{1.75}{0}
\fxrdof{1}{1}{2.5}{0}
\fxrdof{1}{1}{3.5}{0}
\fxrdof{1}{1.5}{0.25}{0}
\fxrdof{1}{1.5}{0.75}{0}
\fxrdof{1}{1.5}{1.25}{0}
\fxrdof{1}{1.5}{1.75}{0}
\fxrdof{1}{2}{0.25}{0}
\fxrdof{1}{2}{0.75}{0}
\fxrdof{1}{2}{1.25}{0}
\fxrdof{1}{2}{1.75}{0}
\fxrdof{1}{2}{2.5}{0}
\fxrdof{1}{2}{3.5}{0}
\fxrdof{1}{3}{0.5}{0}
\fxrdof{1}{3}{1.5}{0}
\fxrdof{1}{3}{2.5}{0}
\fxrdof{1}{3}{3.5}{0}
\fxrlink{2}{0.5}{0.25}{0.5}{0.75}{0}{0}
\fxrlink{2}{0.5}{0.25}{1}{0.25}{0}{0}
\fxrlink{2}{0.5}{0.25}{0}{0.25}{2}{0}
\fxrlink{2}{0.5}{0.75}{0.5}{1.25}{3}{0}
\fxrlink{2}{0.5}{0.75}{1}{0.75}{0}{0}
\fxrlink{2}{0.5}{0.75}{0}{0.75}{2}{0}
\fxrlink{2}{0.5}{1.25}{0.5}{1.75}{3}{0}
\fxrlink{2}{0.5}{1.25}{1}{1.25}{3}{0}
\fxrlink{2}{0.5}{1.25}{0}{1.25}{3}{0}
\fxrlink{2}{0.5}{1.75}{1}{1.75}{3}{0}
\fxrlink{2}{0.5}{1.75}{1}{2.5}{1}{0}
\fxrlink{2}{0.5}{1.75}{0}{1.75}{2}{0}
\fxrlink{2}{0.5}{1.75}{0}{2.5}{2}{0}
\fxrlink{2}{1}{0.25}{1}{0.75}{0}{0}
\fxrlink{2}{1}{0.25}{1.5}{0.25}{0}{0}
\fxrlink{2}{1}{0.75}{1}{1.25}{3}{0}
\fxrlink{2}{1}{0.75}{1.5}{0.75}{0}{0}
\fxrlink{2}{1}{1.25}{1}{1.75}{3}{0}
\fxrlink{2}{1}{1.25}{1.5}{1.25}{3}{0}
\fxrlink{2}{1}{1.75}{1}{2.5}{3}{0}
\fxrlink{2}{1}{1.75}{1.5}{1.75}{3}{0}
\fxrlink{2}{1}{2.5}{1}{3.5}{0}{0}
\fxrlink{2}{1}{2.5}{1.5}{1.75}{3}{0}
\fxrlink{2}{1}{2.5}{2}{2.5}{3}{0}
\fxrlink{2}{1}{2.5}{0}{2.5}{2}{0}
\fxrlink{2}{1}{3.5}{2}{3.5}{0}{0}
\fxrlink{2}{1}{3.5}{0}{3.5}{2}{0}
\fxrlink{2}{1.5}{0.25}{1.5}{0.75}{0}{0}
\fxrlink{2}{1.5}{0.25}{2}{0.25}{1}{0}
\fxrlink{2}{1.5}{0.75}{1.5}{1.25}{3}{0}
\fxrlink{2}{1.5}{0.75}{2}{0.75}{1}{0}
\fxrlink{2}{1.5}{1.25}{1.5}{1.75}{3}{0}
\fxrlink{2}{1.5}{1.25}{2}{1.25}{3}{0}
\fxrlink{2}{1.5}{1.75}{2}{1.75}{3}{0}
\fxrlink{2}{1.5}{1.75}{2}{2.5}{3}{0}
\fxrlink{2}{2}{0.25}{2}{0.75}{1}{1}
\fxrlink{2}{2}{0.25}{2}{0.75}{1}{-1}
\fxrlink{2}{2}{0.25}{2}{1.25}{3}{0}
\fxrlink{2}{2}{0.25}{2}{1.75}{3}{0}
\fxrlink{2}{2}{0.25}{3}{0.5}{1}{0}
\fxrlink{2}{2}{0.75}{2}{1.25}{3}{0}
\fxrlink{2}{2}{0.75}{2}{1.75}{3}{0}
\fxrlink{2}{2}{0.75}{3}{0.5}{1}{0}
\fxrlink{2}{2}{1.25}{2}{1.75}{3}{1}
\fxrlink{2}{2}{1.25}{2}{1.75}{3}{-1}
\fxrlink{2}{2}{1.25}{2}{2.5}{3}{0}
\fxrlink{2}{2}{1.25}{3}{1.5}{3}{0}
\fxrlink{2}{2}{1.75}{2}{2.5}{3}{0}
\fxrlink{2}{2}{1.75}{3}{1.5}{3}{0}
\fxrlink{2}{2}{2.5}{2}{3.5}{3}{0}
\fxrlink{2}{2}{2.5}{3}{2.5}{3}{0}
\fxrlink{2}{2}{3.5}{3}{3.5}{0}{0}
\fxrlink{2}{3}{0.5}{3}{1.5}{3}{0}
\fxrlink{2}{3}{0.5}{4}{0.5}{2}{0}
\fxrlink{2}{3}{1.5}{3}{2.5}{3}{0}
\fxrlink{2}{3}{1.5}{4}{1.5}{3}{0}
\fxrlink{2}{3}{2.5}{3}{3.5}{3}{0}
\fxrlink{2}{3}{2.5}{4}{2.5}{3}{0}
\fxrlink{2}{3}{3.5}{4}{3.5}{2}{0}
\fxrlink{2}{0}{0.25}{0}{0.75}{2}{0}
\fxrlink{2}{0}{0.75}{0}{1.25}{3}{0}
\fxrlink{2}{0}{1.25}{0}{1.75}{3}{0}
\fxrlink{2}{0}{1.75}{0}{2.5}{2}{0}
\fxrlink{2}{0}{2.5}{0}{3.5}{2}{0}
\fxrlink{2}{4}{0.5}{4}{1.5}{2}{0}
\fxrlink{2}{4}{1.5}{4}{2.5}{3}{0}
\fxrlink{2}{4}{2.5}{4}{3.5}{3}{0}
\fxrdof{2}{0.5}{0.25}{0}
\fxrdof{2}{0.5}{0.75}{0}
\fxrdof{2}{0.5}{1.25}{0}
\fxrdof{2}{0.5}{1.75}{0}
\fxrdof{2}{1}{0.25}{0}
\fxrdof{2}{1}{0.75}{0}
\fxrdof{2}{1}{1.25}{0}
\fxrdof{2}{1}{1.75}{0}
\fxrdof{2}{1}{2.5}{0}
\fxrdof{2}{1}{3.5}{0}
\fxrdof{2}{1.5}{0.25}{0}
\fxrdof{2}{1.5}{0.75}{0}
\fxrdof{2}{1.5}{1.25}{0}
\fxrdof{2}{1.5}{1.75}{0}
\fxrdof{2}{2}{0.25}{0}
\fxrdof{2}{2}{0.75}{0}
\fxrdof{2}{2}{1.25}{0}
\fxrdof{2}{2}{1.75}{0}
\fxrdof{2}{2}{2.5}{0}
\fxrdof{2}{2}{3.5}{0}
\fxrdof{2}{3}{0.5}{0}
\fxrdof{2}{3}{1.5}{0}
\fxrdof{2}{3}{2.5}{0}
\fxrdof{2}{3}{3.5}{0}
\fxrdof{2}{0}{0.25}{1}
\fxrdof{2}{0}{0.75}{1}
\fxrdof{2}{0}{1.25}{1}
\fxrdof{2}{0}{1.75}{1}
\fxrdof{2}{0}{2.5}{1}
\fxrdof{2}{0}{3.5}{1}
\fxrdof{2}{4}{0.5}{1}
\fxrdof{2}{4}{1.5}{1}
\fxrdof{2}{4}{2.5}{1}
\fxrdof{2}{4}{3.5}{1}
%
    \end{tikzpicture}}%
  \begin{subfigure}[b]{0.49\textwidth}
    \centering
    \fxrpanel{1}
    \caption{Single-phase convection.}
    \label{fig:momentum:rules:single}
  \end{subfigure}%
  \hfill
  \begin{subfigure}[b]{0.49\textwidth}
    \centering
    \fxrpanel{2}
    \caption{Two-phase transport.}
    \label{fig:momentum:rules:two}
  \end{subfigure}
  \caption[Face states of the dual faces]{Face state of every dual face of
    the component $\velo$ on $\facesaxis{1}$, read from the code's schedules on
    a $4\times4$ mesh of unit cells with walls on all sides whose lower-left
    root is refined once (\code{momentum_rules_data}); every dual face is drawn
    as an edge between the two unknowns it joins (dots; squares: boundary
    unknowns of the phase schedule), arcs where two dual faces of different axes
    join the same pair. Thin grey: regular, the state of \code{scheme};
    \tikz[baseline=-0.6ex]\draw[fxr 1] (0,0) -- (0.45,0);~non-regular or
    incident to a boundary unknown, the state of \code{amr_interface_scheme};
    (a)~boundary dual faces, drawn to their boundary face:
    \tikz[baseline=-0.6ex]\draw[fxr 4] (0,0) -- (0.45,0);~(B1),
    \tikz[baseline=-0.6ex]\draw[fxr 5] (0,0) -- (0.45,0);~(B2), without flux at
    a wall, \tikz[baseline=-0.6ex]\draw[fxr 6] (0,0) -- (0.45,0);~(B3), where the
    level transition meets the wall; (b)~liquid below
    $\coordk{2}=1.1+0.35\coordk{1}$ (blue, PLIC lines), and
    \tikz[baseline=-0.6ex]\draw[fxr 3] (0,0) -- (0.45,0);~CTU: a phase with
    volume at one of the two unknowns has a staggered fraction below
    $\minfracval=\half$ at one of them (here the fractions at the start of a
    sweep, which are the predicted fractions of a vanishing step).}
  \label{fig:momentum:rules}
\end{figure}
